% ============================================================
%  Ejercicio 2.1 — drrA, map2k7 y ndk contra VFDB
% ============================================================
\section{Ejercicio 2.1}

\begin{consigna}
Repita la búsqueda con \texttt{-outfmt 6} y visualice las columnas. Investigue cómo incluir el
alineamiento. Repita para las otras 2 proteínas. ¿Qué puede concluir del análisis?
\end{consigna}

El comando básico con outfmt 6:

\begin{lstlisting}[language=bash]
blastp -query drrA.faa -db VFDB_setB_pro.fas -evalue 1e-5 -outfmt 6
\end{lstlisting}

Las 12 columnas del formato tabular por defecto son:

\begin{center}
\small\ttfamily
qseqid\ \ sseqid\ \ pident\ \ length\ \ mismatch\ \ gapopen\ \ qstart\ \ qend\ \ sstart\ \ send\ \ evalue\ \ bitscore
\end{center}

Para incluir las secuencias alineadas hay que agregar \texttt{qseq} y \texttt{sseq} al final del string de
outfmt:

\begin{lstlisting}[language=bash]
blastp -query drrA.faa -db VFDB_setB_pro.fas -evalue 1e-5 \
  -outfmt "6 qseqid sseqid pident length mismatch gapopen \
           qstart qend sstart send evalue bitscore qseq sseq"
\end{lstlisting}

Resultados para las tres proteínas:

\begin{table}[H]
\centering
\footnotesize
\begin{tabular}{llP{1.8cm}P{2.2cm}}
\toprule
Query & Mejor hit en VFDB & Identidad (\%) & E-value \\
\midrule
drrA   & VFG018236 --- drrA/sidM, \textit{Legionella pneumophila} Ph-1  & 100.0 & 0.0 \\
map2k7 & VFG019987 --- ppkA, \textit{Pseudomonas aeruginosa} LESB58     & 30.2  & $1.6 \times 10^{-18}$ \\
ndk    & VFG031478 --- ndk, \textit{Mycobacterium marinum} M             & 58.8  & $3.2 \times 10^{-52}$ \\
\bottomrule
\end{tabular}
\caption{Mejor hit en VFDB para cada proteína query (blastp, $E < 10^{-5}$).}
\end{table}

\textbf{drrA} tiene un match exacto (100\% de identidad) con un factor de virulencia de VFDB: es la misma
proteína, anotada en VFDB como factor de virulencia de \textit{Legionella pneumophila}. \textbf{ndk} tiene
un hit fuerte (58.8\%, E-value $3.2 \times 10^{-52}$) contra NDK de \textit{Mycobacterium marinum}, que
está anotada como factor de virulencia porque bloquea la maduración del fagosoma; el resultado es
biológicamente coherente. \textbf{map2k7} es el caso más interesante: tiene hits significativos, pero con
solo 30\% de identidad y cubriendo únicamente una región parcial de la proteína (el dominio quinasa).
map2k7 es una quinasa eucariota de la vía de estrés JNK; la similitud refleja conservación del dominio
catalítico quinasa, que está presente en muchas proteínas de organismos muy distintos. No sería correcto
anotarla como factor de virulencia solo por este resultado: el BLAST detecta similitud de dominio, no
equivalencia funcional.

% ============================================================
%  Ejercicio 2.2 — N315 vs VFDB con Bio.SearchIO
% ============================================================
\section{Ejercicio 2.2}

\begin{consigna}
Buscar todos los factores de virulencia de VFDB sobre el proteoma de \textit{S.\ aureus} N315.
Generar una tabla con: query, hit, cobertura de la query, cobertura del hit, identidad y especie del hit.
\end{consigna}

Corrimos BLAST del proteoma completo de N315 contra VFDB con el siguiente comando (incluimos \texttt{qlen}
y \texttt{slen} para poder calcular las coberturas en el script):

\begin{lstlisting}[language=bash]
blastp -query staph_aureus_N315_proteome.faa -db VFDB_setB_pro.fas \
  -evalue 1e-5 -num_threads 4 -max_target_seqs 10 \
  -outfmt "6 qseqid sseqid pident length mismatch gapopen \
           qstart qend sstart send evalue bitscore qlen slen \
           qseq sseq stitle" \
  -out N315_vs_VFDB_blast.tsv
\end{lstlisting}

Para analizar la salida usamos \texttt{Bio.SearchIO} de Biopython
(\texttt{scripts/ejercicio\_2\_2\_biopython.py}):

\begin{lstlisting}[language=Python]
import Bio.SearchIO as bpio

campos = ["qseqid","sseqid","pident","length","mismatch","gapopen",
          "qstart","qend","sstart","send","evalue","bitscore",
          "qlen","slen","stitle"]

resultados = list(bpio.parse("N315_vs_VFDB_blast.tsv",
                              "blast-tab", fields=campos))

for query_result in resultados:
    for hit in query_result:
        for hsp in hit:
            cob_query = (hsp.query_end - hsp.query_start) * 100 \
                        / query_result.seq_len
            cob_hit   = (hsp.hit_end - hsp.hit_start) * 100 \
                        / hit.seq_len
            especie = hit.description  # extraído de stitle
            # filtrar por cobertura del hit >= 70% y armar la tabla
\end{lstlisting}

Con el criterio de cobertura del hit $\geq 70\%$ identificamos \textbf{503 proteínas} del proteoma de N315
con similitud significativa a factores de virulencia de VFDB. La tabla completa está en
\texttt{resultados/N315\_vs\_VFDB\_biopython.tsv} con columnas: query, hit, cobertura\_query (\%),
cobertura\_hit (\%), identidad (\%), evalue, bitscore, especie\_hit.

% ============================================================
%  Ejercicio 2.3 — Enterotoxinas en N315
% ============================================================
\section{Ejercicio 2.3}

\begin{consigna}
\textbf{a)} BLAST del proteoma de N315 contra la base de enterotoxinas. ¿Cuáles están presentes? \\
\textbf{b)} Programar un criterio de decisión por identidad/cobertura. Si más de un hit cumple las
expectativas, indicar que se encontró más de una.
\end{consigna}

Usamos el siguiente criterio de decisión: identidad $\geq 80\%$, cobertura de la query $\geq 80\%$ y
cobertura del hit $\geq 80\%$. El script (\texttt{scripts/ejercicio\_2\_3\_biopython.py}) lee la salida
tabular, aplica estos tres filtros y para cada proteína query lista las enterotoxinas que los pasan. Si
más de una enterotoxina cumple para la misma proteína, la decisión queda como \texttt{multiple}.

\medskip
Enterotoxinas detectadas en el proteoma de \textit{S.\ aureus} N315:

\begin{table}[H]
\centering
\footnotesize
\begin{tabular}{llP{1.5cm}P{1.8cm}P{2.2cm}}
\toprule
Proteína N315   & Enterotoxina(s)  & Decisión & Identidad (\%) & E-value \\
\midrule
WP\_000034846.1 & sep              & única    & 99.6  & $\sim 0$ \\
WP\_000278088.1 & sec3; sec2; sec1 & multiple & 100.0 & $\sim 0$ \\
WP\_000475325.1 & selx             & única    & 96.1  & $8.6 \times 10^{-146}$ \\
WP\_000713847.1 & sei              & única    & 100.0 & $1.5 \times 10^{-180}$ \\
WP\_000736712.1 & seg              & única    & 100.0 & $\sim 0$ \\
WP\_000746599.1 & sel              & única    & 99.2  & $2.3 \times 10^{-176}$ \\
WP\_000821658.1 & sem; selv        & multiple & 92.5  & $6.0 \times 10^{-167}$ \\
WP\_001035599.1 & tsst-1           & única    & 99.1  & $3.6 \times 10^{-174}$ \\
WP\_001236362.1 & sen              & única    & 100.0 & $\sim 0$ \\
WP\_010922839.1 & seo              & única    & 99.6  & $\sim 0$ \\
\bottomrule
\end{tabular}
\caption{Enterotoxinas estafilocócicas detectadas en \textit{S.\ aureus} N315
(criterio: identidad $\geq 80\%$, coberturas $\geq 80\%$).}
\end{table}

Encontramos 10 enterotoxinas en total. Los dos casos \texttt{multiple} tienen una explicación biológica
clara: sec1, sec2 y sec3 son variantes muy parecidas de la enterotoxina estafilocócica tipo C, con
secuencias tan similares entre sí que con nuestro criterio de 80\% de identidad no es posible
distinguirlas; lo mismo pasa con sem y selv. En esos casos preferimos reportar \texttt{multiple} en lugar
de elegir una al azar, ya que forzar una clasificación única sería incorrecto.
